\chapter{Implementation}
\label{ch:implementation}

\section{Ingestion}
\label{sec:impl-ingestion}

\subsection{Partitioning Is a Correctness Decision}
\label{sec:impl-partitioning}

The topic \code{fleet.telemetry} has 6 partitions, \textbf{keyed by
\code{vehicle\_id}}. Throughput has nothing to do with it. Kafka
guarantees ordering only \emph{within} a
partition~\cite{kreps2011kafka}, and the speed layer's per-vehicle state
machine, which tracks when each vehicle last went idle, assumes it sees
one vehicle's events in the order they happened. Keying by
\code{vehicle\_id} is what gives it that. Key round-robin instead and
every idle measurement quietly goes wrong, without anything anywhere
raising an error. The dead-letter topic \code{fleet.telemetry.dlq} has one
partition, since it needs no ordering and one partition keeps every
rejected event in a single readable stream.

\subsection{Fault Injection}
\label{sec:impl-faults}

The simulator corrupts its own output, so we can show the error handling
working rather than claim it does. Table~\ref{tab:faults} gives
configured against measured rates over roughly 2,700 events. Validation
rejects rather than repairs: coerce a malformed field and you
get a confidently wrong money figure downstream, which is worse than a
rejection somebody can see. The DLQ carries the reason code \emph{and the
original payload}, so once a producer is fixed its rejected events can be
replayed. The observed distribution was \code{unparseable\_json}~58,
\code{coords\_out\_of\_bounds}~23, \code{bad\_status}~17,
\code{negative\_fare}~12 and \code{negative\_speed}~10.

\begin{table}[htbp]
\centering
\caption{Injected fault rates, configured against measured, and what each
fault demonstrates.}
\label{tab:faults}
\begin{tabularx}{\textwidth}{@{}llrX@{}}
\toprule
\textbf{Fault} & \textbf{Config.} & \textbf{Measured} &
\textbf{What it proves}\\
\midrule
Malformed JSON & 0.5\% & 0.40\% &
  The archiver preserves unparseable bytes; the speed layer
  dead-letters rather than crashing\\
Schema-invalid & 0.5\% & 0.77\% &
  \code{common/validation.py} rejects identically in both layers\\
Duplicate \code{event\_id} & 1.0\% & 1.03\% &
  Deduplication works; both copies are \emph{valid}, so nothing else
  can stop them\\
Late by 1--20 sim min & 2.0\% & 2.09\% &
  The consistency argument; see Section~\ref{sec:consistency}\\
\bottomrule
\end{tabularx}
\end{table}

A second simulator drops one expense CSV per simulated day. Files are
versioned, so a partner correcting an earlier submission drops a
\code{v2} file for a date already reconciled. That is the event the whole
architecture exists to absorb (Section~\ref{sec:recompute}). The delivery
SLA governs first delivery only. An early version checked it
against every version, which meant a corrected \code{v2} marked every
resubmission late.

\section{The Speed Layer}
\label{sec:impl-speed}

\subsection{The Archiver Does as Little as Possible}
\label{sec:impl-archiver}

The archiver is a \emph{separate Spark application} with its own
checkpoint and \textbf{no business logic at all}. Every past day is
recomputed from the master dataset, so if the master dataset is corrupt
then all history is wrong and nothing can recover it. Keeping the
archiver separate means a bug in the speed layer's validation, windowing
or state handling \emph{cannot reach it}. The worst such a bug can do is
produce wrong dashboards, and the next batch run overwrites those.
Unparseable records are archived under \code{sim\_date=unknown} with
their raw bytes intact. An archive that quietly dropped them would make
``how much bad data did we receive?'' unanswerable forever.

\subsection{Stateful Idle Alerts}
\label{sec:impl-idle}

No window or aggregation can answer ``has this vehicle been idle for 45
simulated minutes?''. Answering it needs memory of \emph{when} the
vehicle last stopped being idle, carried across micro-batches, plus the
ability to fire when no new event arrives at all, since a silent
vehicle is exactly the case we care about.
\code{applyInPandasWithState} is the only Structured Streaming construct
that offers both~\cite{spark_ss_guide}. Four details decide whether it
works, and each is a defect if you get it wrong:

\begin{enumerate}
  \item \textbf{The stopwatch starts on the \emph{transition} into
        idle}, not on every idle ping. Resetting per ping means the timer
        never advances and no alert can fire.
  \item \textbf{Events are sorted by event time before folding}, because
        a micro-batch may interleave partitions.
  \item \textbf{A stale out-of-order event does not rewind the state.}
        The event is not lost, since the archive holds it and the batch
        layer counts it, but rewinding would corrupt the idle
        measurement.
  \item \textbf{\code{EventTimeTimeout}, not
        \code{ProcessingTimeTimeout}.} A processing-time timeout fires
        after $N$ \emph{real} seconds, so changing the clock compression
        would silently change the business rule.
\end{enumerate}

This design has a consequence we would rather state than hide. An
event-time timeout fires only when the watermark advances, and the
watermark advances only when other events
arrive~\cite{akidau2015dataflow}. If the whole fleet went silent at once,
no timeout would fire at all. A different control covers that case: the
\code{TelemetryNotProduced} alert, which watches the producer directly.

\subsection{Idempotent Sinks}
\label{sec:impl-sinks}

\code{foreachBatch} gives \textbf{at-least-once} delivery. If the driver
dies after writing a micro-batch but before committing offsets, Spark
re-runs the batch and the sink sees the same rows
twice~\cite{armbrust2018structured}. With plain inserts, every dashboard
figure would ratchet upward on each restart and never come back down.
Every sink therefore upserts on a natural key:
\code{rt\_vehicle\_state} on \code{(vehicle\_id)},
\code{rt\_zone\_hourly} on \code{(zone\_id, window\_start)},
\code{rt\_vehicle\_daily} on \code{(vehicle\_id, sim\_date)}, and
\code{idle\_alerts} on \code{(vehicle\_id, opened\_at)}. A repeated write
then sets the same values it set the first
time~\cite{kleppmann2017ddia}. It also lets the windowed aggregations use
\code{update} output mode, which they need: under \code{append}, a window
is withheld until the watermark passes its end, so the \emph{current}
simulated hour, the one thing a live dashboard exists to show, would
never appear.

\section{The Batch Layer and Orchestration}
\label{sec:impl-batch}

\subsection{The Batch Layer Trusts Nothing}
\label{sec:impl-trust}

\code{compute\_vehicle\_day} re-reads the raw archive and re-applies the
same validation and deduplication from scratch. It never reads the speed
layer's output, and it applies \textbf{no watermark}. It is reading a
complete, closed day, so lateness means nothing to it, and it
deduplicates across the whole day instead of within a watermark window.
That independence is the only reason the drift measurement in
Section~\ref{sec:consistency} means anything: the two figures really are
arrived at twice. Its outputs per vehicle-day are \code{trips},
\code{revenue}, \code{gps\_km} (haversine between consecutive pings,
discarding jumps above 150\,km/h), \code{online\_min},
\code{on\_trip\_min}, \code{idle\_min} and \code{utilization}.

The reconciliation joins telemetry to expenses with a full outer join,
and the choice is not accidental. An inner join would quietly drop
the two findings the business most needs to see: telemetry with no
expense row (\code{missing\_costs}, meaning we are blind to a vehicle's
costs) and an expense row with no telemetry
(\code{missing\_telemetry}, meaning we are being billed for a vehicle
that never reported).

\subsection{Schedule and Recomputation}

Airflow's scheduler lives in real time and knows nothing about the
compressed clock. So the DAG runs every 2 real minutes and its
\emph{first} task works out which simulated day is pending. The schedule
is a poll; the date selection comes from the data. We did try expressing
the schedule in simulated time, and got an interval 60 times longer than
we meant. The DAG has nine tasks: \code{find\_pending\_date},
\code{check\_expense\_file}, \code{validate\_expenses},
\code{compute\_vehicle\_day}, \code{reconcile\_profitability},
\code{load\_batch\_views}, \code{compute\_drift},
\code{generate\_report} and \code{record\_run}.

\code{load\_batch\_views} deletes every row for the target
\code{sim\_date} and inserts the recomputed set inside a single
transaction, so rerunning a date produces the same row count with updated
values and a partially applied recompute is impossible. This is the
property the architecture exists to provide, demonstrated in
Section~\ref{sec:recompute}.

\section{Storage and Serving}
\label{sec:impl-serving}

Both layers' outputs live in one PostgreSQL instance, every \code{rt\_*}
table keyed on its natural key so the streaming upserts work
(Appendix~\ref{app:schema}). The API exposes eleven documented endpoints
(Appendix~\ref{app:api}), the consolidated daily report is served at
\code{GET /reports/\{sim\_date\}}, and two Grafana dashboards read the
same tables.

\subsection{The Lambda Merge Is Explicit, Never Blended}
\label{sec:impl-merge}

\code{GET /vehicles/\{id\}} returns three blocks, each labelled with its
own \code{source} and \code{as\_of}: \code{live} and \code{today} carry
\code{"source": "speed"}, while \code{history} carries
\code{"source": "batch"}. We do not merge them into a single
figure. The two layers answer different questions under different
guarantees, and a consumer has to be able to see at a glance which
numbers are safe to quote to finance. Blending them is how a dashboard
estimate ends up in a financial statement. ``Now'' is always
\emph{simulated} now, taken from \code{max(last\_event\_time)} rather
than the database clock, so if ingestion stalls then ``now'' stops
advancing and the active-vehicle count drops to zero instead of quietly
lying.

\section{Observability}
\label{sec:impl-observability}

Four questions drive every metric. Nothing is collected merely because it
was easy to collect.

\emph{Is data arriving?} \code{fleet\_events\_produced\_total} and
\code{fleet\_faults\_injected\_total} answer that, and between them they
separate ``the producer is silent'' from ``the producer is gone'', which
is why the alert tests both \code{increase()==0} and \code{absent()}.
\emph{Is it being processed?}
\code{fleet\_stream\_last\_progress\_timestamp} and its siblings, all
recorded per query rather than per process, because a container can look
perfectly healthy while one of its five queries sits frozen. \emph{Is it
correct?} \code{fleet\_dlq\_events\_total},
\code{fleet\_late\_rows\_dropped\_total} and
\code{fleet\_speed\_batch\_drift\_ratio}, which measure imprecision we
know about and have accepted, on the principle that a pipeline unable to
quantify its own error has no business handling money. \emph{Is the
business healthy?} \code{fleet\_open\_idle\_alerts} and the profitability
flags, kept well away from pipeline health, because
\code{IdleVehiclesHigh} fires when everything works and the fleet simply
is not earning.

Every log line is one JSON object carrying \code{ts}, \code{level},
\code{service}, \code{stage}
(\code{ingestion}/\code{processing}/\code{storage}/\code{serving}/\code{orchestration}),
\code{event}, \code{msg} and \code{sim\_time}. That last field
matters most here, because on a compressed clock a wall-clock timestamp
tells you nothing about \emph{which simulated hour} a line belongs to.
Per-event logging is sampled 1-in-$N$: at 50 vehicles emitting every 2
seconds, an INFO line per event buries everything else, so the useful
granularity is one line per micro-batch carrying counts.

We do not implement distributed tracing at all. What we implement is
\textbf{correlation-ID propagation}. \code{event\_id} travels producer
$\rightarrow$ Kafka $\rightarrow$ archiver $\rightarrow$ DLQ
$\rightarrow$ batch dedup. \code{vehicle\_id} travels through every stage
and doubles as the partition key. \code{run\_id} tags every Airflow task
and batch table, and \code{sim\_date} identifies the lake partition and
the report filename. So a rejected event can be traced from its DLQ
reason back to its original payload, and a reconciled figure back to the
run and expense-file version behind it. What we give up is per-request
span timing across service boundaries. In a pipeline whose hops are Kafka
offsets and Parquet partitions rather than RPC calls, spans would add
infrastructure without answering a question we actually have.
Section~\ref{sec:production} says when that stops being true.

\subsection{Alerting Decisions Worth Defending}
\label{sec:impl-alerting}

\textbf{Kafka lag comes from Spark's own query progress, not from a lag
exporter.} Structured Streaming never commits consumer-group offsets. It
tracks them in its checkpoint instead, which means
\code{kafka-consumer-groups.sh} and every off-the-shelf exporter report
\emph{nothing} for these consumers. The number exists in one place only:
the progress event's source metrics.

\textbf{Batch metrics are exported by the API from PostgreSQL, not
pushed.} Airflow tasks are short-lived and Prometheus pulls, so by the
time a scrape happens the task has exited. A Pushgateway would work, but
it is one more service holding job state indefinitely with no notion of
staleness, and the Prometheus documentation discourages it for exactly
this pattern~\cite{prometheus_docs}. The batch layer already persists its
outcomes because the API and the report need them, so re-exporting costs
us nothing and means the metric and the report cannot disagree.

Ten alert rules are defined (Appendix~\ref{app:alerts}), each with a
\code{for:} duration, a severity and a runbook entry.
Section~\ref{sec:alerting-results} reports which have been verified
firing and delivered.
